% !TEX root = ../main.tex
%
\section{Materials}
\label{sec:materials}
%TODO: expand on this

\subsection{Software}
\label{sec:materials:software}

We introduce an open-source, lightweight, purpose-built framework for generating synthetic discussions.\footnote{\librarylink} The key features of the framework are: 
\begin{itemize}[nosep,noitemsep]
	\item Three core functionalities: generating discussion setups (selecting participants, topics, roles, etc.), executing, and annotating them according to user-provided parameters.\footnote{Although we did not use LLM annotation in this study.}
	\item Built-in fault tolerance (automated recovery and intermittent saving) and file logging to support experiments running for an extended period of time.
	\item Applicability to other domains and use cases such as debate simulation, or discussion-based agentic systems through an adaptable, plug-and-play, object-oriented library interface.
	\item Availability via \texttt{pip}.
\end{itemize}
Our library is released under the GNU General Public License v3 (GPLv3). We also release the code used for running the experiments and analyzing, visualizing and exporting the results under the same license.\footnote{\experimentslink}
%TODO: add links


\subsection{Datasets}
\label{sec:materials:data}

\begin{figure}[ht]
	\centering
	\includegraphics[width=0.45\columnwidth]{toxicity_distribution_human_vs_synthetic.png}
	\includegraphics[width=0.45\columnwidth]{comment_len_model.png}
	\Description{Distribution plot showing toxicity ratings between human and synthetic discussions. Distribution plot showing comment length in words by model, including human discussions. }
	\caption{\textbf{Left}: Distribution of toxicity ratings between the main synthetic dataset and the human seed-opinion dataset. We remove empty comments. \textbf{Right}: Comment length of human and LLM discussions by model. Qwen-32B is the only model that occasionally chooses not to speak-this is largely due to it being the only facilitator model to occasionally not intervene (\S\ref{sec:facilitation:interventions}). Synthetic discussions closely approximate the means of human toxicity and \diversity distributions with less variance. }
	\label{fig:human-llm-comparison}
\end{figure}


We present a dataset designed to support research in human-computer interaction, social dynamics, and language modeling for facilitation tasks. Figure~\ref{fig:human-llm-comparison} shows that the main synthetic dataset exhibits toxicity levels and comment lengths that broadly align with those of the human seed-opinion dataset. However, in both cases the synthetic datasets display noticeably sharper distributions centered around the human mean. This pattern confirms studies showing that LLMs can replicate common human behaviors and opinions, but often fail to capture fringe, long-tail behaviors~\cite{wu2025llm, madden2025evaluating,ivetta2025synthetic, kaiser2025simulating}.

Our dataset is primarily intended to serve as a resource for analyzing off-the-shelf LLMs in the task of online facilitation, as well as the behaviors and language used by LLM participants since such datasets are extremely limited in literature~\cite{korre2025evaluation}. Additionally, it could be used for fine-tuning LLMs for facilitation, as recent work demonstrates that task-specific LLMs can be finetuned with data from their own simulations with adequate filtering~\cite{ulmer2024}--which could be performed by our evaluation pipeline (\S\ref{sec:methodology:evaluation}). However, it may require filtering to ensure the quality and representativeness of the discussions~\cite{ulmer2024}--in our case filtering out discussions with constant facilitator interventions or low/extremely high diversity. This restriction can be overcome by scaling up the experiments, which is feasible due to their very low computational cost (\S\ref{sec:scalability:cost}), and by exclusively choosing appropriate models, such as LLaMa-8B and Qwen-8B (\S\ref{sec:scalability:model}).


The supplementary ablation dataset, as well as the code for the analysis and the graphs present in this paper, can be found in the project repository. The dataset is licensed under a CC BY-SA license. We include relevant statistics for the main, ablation and human discussion~\cite{chang-danescu-niculescu-mizil-2019-trouble} datasets in Table~\ref{tab:dataset-stats}.

\begin{table}[ht]
	\caption{Dataset statistics for the main dataset, and the ablation experiments.}
	\label{tab:dataset-stats}
	\begin{tabular}{lrr}
		\toprule
		& Main & Ablation  \\
		\midrule
		Mean Comments per Discussion & 26.25 & 30.00\\
		Mean Words per Comment & 53.52 & 48.18 \\
		Std. Comments per Discussion & 6.50 & 0.00  \\
		Std. Words per Comment & 46.56 & 34.66  \\
		\# Comments & $18{,}900$ & $32{,}400$  \\
		\# Discussions & 720 & 1080 \\
		\bottomrule
	\end{tabular}
\end{table}
